\subsection{Método de Shubert-Piyavskii}

% ---------------------------------------------------------------
% 1. Explicación conceptual y Placeholder gráfico
% ---------------------------------------------------------------
\begin{frame}{Shubert-Piyavskii: Explicación Conceptual}
    \begin{block}{Concepto principal}
        \footnotesize
        El método de Shubert-Piyavskii es un algoritmo de optimización global para funciones 
        que cumplen con la condición de Lipschitz (su tasa de cambio está acotada por una constante $\ell$). 
        A diferencia de los métodos de búsqueda local, este construye iterativamente una 
        \textbf{función de límite inferior} que aproxima la función real desde abajo.
    \end{block}

    \scriptsize
    \begin{itemize}
        \item En cada paso, evalúa la función donde el límite inferior es \textbf{mínimo}.
        \item Garantiza encontrar el óptimo global aislando regiones prometedoras.
    \end{itemize}
\end{frame}

% ---------------------------------------------------------------
% 2. Explicación Visual de la Constante de Lipschitz
% ---------------------------------------------------------------
\begin{frame}{Shubert-Piyavskii: La Constante de Lipschitz}
    \begin{block}{Interpretación Geométrica}
        \footnotesize
        La condición \(|f(x) - f(y)| \le \ell|x - y|\) significa que la tasa de cambio de la función 
        está acotada por \(\ell\).
    \end{block}

   \begin{center}
    \begin{tikzpicture}[x=1.1cm, y=0.6cm, font=\scriptsize]
        
        % Sombreado de la región válida (Cono L=1) restringido a las abscisas x e y
        \fill[VerdeNeon, opacity=0.15] (2,1) -- (4,3) -- (2,5) -- cycle;
        \fill[VerdeNeon, opacity=0.15] (6,1) -- (4,3) -- (6,5) -- cycle;

        % Ejes (Diferenciados de las abscisas "x" e "y")
        \draw[thick, GrisTexto, -stealth] (0,0) -- (8.5,0) node[right] {};
        \draw[thick, GrisTexto, -stealth] (0,0) -- (0,6.5) node[above] {};

        % Abscisas "x" e "y" solicitadas
        \draw[GrisTexto, thick] (2, 0.1) -- (2, -0.1) node[below, font=\small] {\(x\)};
        \draw[GrisTexto, thick] (6, 0.1) -- (6, -0.1) node[below, font=\small] {\(y\)};
        
        % Líneas punteadas verticales desde x e y hasta arriba de las líneas verdes con sus etiquetas
        \draw[GrisTexto, dashed] (2,0) -- (2,5.7);
        \node[above, text=GrisTexto] at (2,5.7) {\(f(x)\)};

        \draw[GrisTexto, dashed] (6,0) -- (6,5.7);
        \node[above, text=GrisTexto] at (6,5.7) {\(f(y)\)};
        
        % Texto explicativo de la constante
        \node[GrisTexto, font=\footnotesize\bfseries, fill=FondoOscuro, rounded corners, inner sep=3pt] at (4, 5.8) {\(\ell = 1\)};

        % Trazado de las líneas intersectando en (4,3)
        % Las líneas se dibujan de izquierda a derecha para que 'pos=0.85' quede a la derecha
        
        % m = 2 (Excede L, rojo punteado)
        \draw[RojoAlerta, thick, dashed] (2.5,0) -- (5.5,6) 
            node[sloped, above, pos=0.85, text=RojoAlerta] {\(m=2\)};

        % m = 1 (Límite superior, verde)
        \draw[VerdeNeon, thick] (1,0) -- (7,6) 
            node[sloped, above, pos=0.85, text=VerdeNeon] {\(m=1\)};

        % m = 1/2 (Válido, verde)
        \draw[VerdeNeon, thick] (0.5,1.25) -- (7.5,4.75) 
            node[sloped, above, pos=0.85, text=VerdeNeon] {\(m=1/2\)};

        % m = 0 (Válido, verde)
        \draw[VerdeNeon, thick] (0.5,3) -- (7.5,3) 
            node[sloped, above, pos=0.85, text=VerdeNeon] {\(m=0\)};

        % m = -1/2 (Válido, verde)
        \draw[VerdeNeon, thick] (0.5,4.75) -- (7.5,1.25) 
            node[sloped, above, pos=0.85, text=VerdeNeon] {\(m=-1/2\)};

        % m = -1 (Límite inferior, verde)
        \draw[VerdeNeon, thick] (1,6) -- (7,0) 
            node[sloped, above, pos=0.85, text=VerdeNeon] {\(m=-1\)};

        % m = -2 (Excede -L, rojo punteado)
        \draw[RojoAlerta, thick, dashed] (2.5,6) -- (5.5,0) 
            node[sloped, above, pos=0.85, text=RojoAlerta] {\(m=-2\)};

        % Punto de intersección central
        \filldraw[CelesteBrillante] (4,3) circle (2pt);

    \end{tikzpicture}
\end{center}
\end{frame}

% ---------------------------------------------------------------
% 3. Construcción del cono inicial
% ---------------------------------------------------------------
\begin{frame}{Shubert-Piyavskii: Construcción Inicial}
    \begin{block}{Construcción del cono}
        \footnotesize
        La base del método consiste en evaluar un punto inicial $x_0$. Dado que sabemos que la 
        función no puede descender más rápido que la constante de Lipschitz $\ell$, trazamos dos líneas 
        descendentes de pendiente $\pm \ell$ a partir del punto evaluado $(x_0, f(x_0))$.
    \end{block}

    \begin{center}
    \begin{tikzpicture}[x=0.9cm, y=0.6cm, font=\scriptsize]
        % Caja exterior (ejes)
        \draw[thick, GrisTexto] (0,0) rectangle (10,5.5);
        
        % Ticks inferiores y etiquetas del eje X
        \draw[GrisTexto] (5.0,0.1) -- (5.0,-0.1) node[below] {$x_{0}$};
        
        % Ticks superiores (referencia visual)
        \draw[GrisTexto] (5.0,5.5) -- (5.0,5.4);
        
        % Etiqueta del eje Y
        \node[GrisTexto, rotate=90] at (-0.3, 2.75) {$y$};
        
        % Leyenda en la parte superior (centrada y con colores)
        \draw[white, thick] (2.5, 5.0) -- (3.2, 5.0) node[right, text=white] {$f(x)$};
        \draw[NaranjaBase, thick] (5.5, 5.0) -- (6.2, 5.0) node[right, text=NaranjaBase] {límite inferior};
        
        % Función f(x) (Color blanco, ajustada para no cruzar el límite inferior)
        \draw[white, thick] plot [smooth, tension=0.7] coordinates {
            (1.0, 4.8) 
            (3.0, 3.5) 
            (5.0, 4.0) 
            (7.0, 4.6) 
            (9.0, 3.8)
        };
        
        % Cono de Límite inferior (NaranjaBase, vértice exactamente en x=5)
        \draw[NaranjaBase, thick] (1.0, 0.8) -- (5.0, 4.0) -- (9.0, 0.8);
        
        % Triángulo de pendiente (l) calculado sobre la recta izquierda
        \draw[GrisTexto] (2.0, 1.6) -- (3.0, 1.6) node[midway, below] {$1$};
        \draw[GrisTexto] (3.0, 1.6) -- (3.0, 2.4) node[midway, right] {$\ell$};
        
        % Punto evaluado x_0 centrado
        \filldraw[CelesteBrillante] (5.0, 4.0) circle (1.5pt) 
            node[above, font=\footnotesize] {$(x_0, f(x_0))$};

    \end{tikzpicture}
    \end{center}
    
    \footnotesize
    Estas dos líneas forman un  \textbf{cono en V invertida} debajo del punto. 
    La función real jamás cruzará por debajo del cono, nuestro primer límite inferior global.
\end{frame}

% ---------------------------------------------------------------
% 4. Detalle matemático de las líneas
% ---------------------------------------------------------------
\begin{frame}{Shubert-Piyavskii: Detalle Matemático}
    \begin{alertblock}{La Ecuación del Límite Inferior}
        \small
        Matemáticamente, la condición de Lipschitz establece que para cualquier par de puntos $x$ e $x_i$:
        $|f(x) - f(x_i)| \le \ell |x - x_i|$
    \end{alertblock}

    \footnotesize
    Despejando para encontrar el límite inferior respecto a un punto ya evaluado $x_i$, obtenemos 
    la ecuación de las líneas rectas del cono:
    
    \begin{equation*}
        f(x) \ge f(x_i) - \ell|x - x_i|
    \end{equation*}

    \scriptsize
    \begin{itemize}
        \item Si $x > x_i$ (a la derecha), la ecuación de la recta es $y = f(x_i) - \ell(x - x_i)$ (pendiente $-\ell$).
        \item Si $x < x_i$ (a la izquierda), la ecuación de la recta es $y = f(x_i) + \ell(x - x_i)$ (pendiente $+\ell$).
    \end{itemize}
    
    \vspace{0.5em}
    \footnotesize
    Con múltiples puntos evaluados, el límite inferior global $F(x)$ se define como el máximo de todos los conos superpuestos:
    $F(x) = \max_i \{ f(x_i) - \ell|x - x_i| \}$
\end{frame}

% ---------------------------------------------------------------
% 5. Explicación de Iteración (Visual)
% ---------------------------------------------------------------
\begin{frame}{Shubert-Piyavskii: Desarrollo de la iteración}
    \begin{exampleblock}{Iteración}
        \footnotesize
        La intersección de múltiples conos invertidos de iteraciones pasadas forma un perfil aserrado (límite inferior). Se toma el vértice más bajo del perfil aserrado y se evalúa la función arriba para trazar el siguiente cono que levanta esta parte del perfil.
    \end{exampleblock}

   \begin{center}
\begin{tikzpicture}[x=0.9cm, y=0.7cm, font=\scriptsize]
    % Caja exterior (ejes)
    \draw[thick, GrisTexto] (0,0) rectangle (10,5.5);
    
    % Ticks inferiores y etiquetas del eje X
    \draw[GrisTexto] (1.0,0.1) -- (1.0,-0.1) node[below] {\(a = x^{(2)}\)};
    \draw[GrisTexto] (4.6,0.1) -- (4.6,-0.1) node[below] {\(x^{(1)}\)};
    \draw[GrisTexto] (7.0,0.1) -- (7.0,-0.1) node[below] {\(x^{(4)}\)};
    \draw[GrisTexto] (9.0,0.1) -- (9.0,-0.1) node[below] {\(b = x^{(3)}\)};
    \node[GrisTexto] at (5.0, -0.8) {\(x\)};
    
    % Ticks superiores (referencia visual)
    \draw[GrisTexto] (1.0,5.5) -- (1.0,5.4);
    \draw[GrisTexto] (4.6,5.5) -- (4.6,5.4);
    \draw[GrisTexto] (7.0,5.5) -- (7.0,5.4);
    \draw[GrisTexto] (9.0,5.5) -- (9.0,5.4);
    
    % Etiqueta del eje Y
    \node[GrisTexto, rotate=90] at (-0.3, 2.75) {\(y\)};
    
    % Leyenda en la parte superior (actualizada al color naranja)
    \draw[white, thick] (2.5, 5.0) -- (3.2, 5.0) node[right, text=white] {\(f(x)\)};
    \draw[NaranjaBase, thin] (5.5, 5.0) -- (6.2, 5.0) node[right, text=NaranjaBase] {límite inferior};
    
    % Función f(x) (Color blanco, ajustada para NUNCA cruzar el límite inferior)
    \draw[white, thick] plot [smooth, tension=0.5] coordinates {
        (1.0, 4.8) 
        (2.8, 3.2) 
        (4.6, 4.2) 
        (5.6, 4.8) 
        (7.0, 2.8) 
        (8.2, 2.6) 
        (9.0, 3.6)
    };
    
    % Límite inferior anterior (Dientes de sierra en NaranjaBase y línea delgada)
    \draw[NaranjaBase, thin] (1.0, 4.8) -- (3.0, 1.8) -- (4.6, 4.2) -- (7.0, 0.6) -- (9.0, 3.6);
    
    % Proyección vertical desde el punto mínimo hacia la función (Punteada blanca)
    \draw[-stealth, dashed, white, thick] (7.0, 0.7) -- (7.0, 2.65);
    
    % Punto inferior más bajo (Resaltado en NaranjaBase)
    \filldraw[NaranjaBase] (7.0, 0.6) circle (1.5pt);
    
    % Nuevo cono evaluado en x^(4) (CelesteBrillante)
    \draw[CelesteBrillante, thick] (6.27, 1.7) -- (7.0, 2.8) -- (7.73, 1.7);
    
    % Puntos evaluados reales f(x) (CelesteBrillante)
    \filldraw[CelesteBrillante] (1.0, 4.8) circle (1.5pt); % x^(2)
    \filldraw[CelesteBrillante] (4.6, 4.2) circle (1.5pt); % x^(1)
    \filldraw[CelesteBrillante] (9.0, 3.6) circle (1.5pt); % x^(3)
    \filldraw[CelesteBrillante] (7.0, 2.8) circle (1.5pt); % x^(4)
\end{tikzpicture}
\end{center}
\end{frame}

% ---------------------------------------------------------------
% 6. Pseudocódigo
% ---------------------------------------------------------------
\begin{frame}{Shubert-Piyavskii: Pseudocódigo}
    \scriptsize
    \begin{algorithmic}[1]
        \Require \(f\) Lipschitz con constante \(\ell > 0\), dominio \([a,b]\), tolerancia \(\epsilon\)
        \State \(x_0 \gets (a+b)/2, \quad y_0 \gets f(x_0)\)
        \State \(V \gets \{\text{intersección del cono en } x_0 \text{ con } x=a \text{ y } x=b\}\) \Comment{Vértices iniciales}
        \While{\textbf{true}}
            \State \((x^*, y^*) \gets \arg\min_{(x,y) \in V} y\) \Comment{Vértice inferior actual}
            \State \(y_{\text{real}} \gets f(x^*)\)
            \If{\(y_{\text{real}} - y^* < \epsilon\)} \Comment{Criterio de convergencia}
                \State \Return \(x^*\)
            \EndIf
            \State \(V \gets V \setminus \{(x^*, y^*)\}\) \Comment{Remover punto evaluado}
            \State \((x_L, y_L) \gets \text{intersección izquierda con el segmento previo}\)
            \State \((x_R, y_R) \gets \text{intersección derecha con el segmento posterior}\)
            \State \(V \gets V \cup \{(x_L, y_L), \, (x_R, y_R)\}\)
        \EndWhile
    \end{algorithmic}
    \vspace{0.4em}
    {\tiny \textit{Nota:} El método requiere conocer o sobreestimar adecuadamente la constante de Lipschitz \(\ell\).}
\end{frame}

\begin{frame}{Shubert-Piyavskii: Planteamiento}
    \textbf{Problema:} Minimizar \(f(x) = (x - 0.5)^2\) en \([a, b] = [-1, 1]\).
    
    \vspace{0.3cm}
    \textbf{1. Constante de Lipschitz (\(l\)):}
    \begin{itemize}
        \item Derivada: \(f'(x) = 2(x - 0.5)\).
        \item Máximo valor absoluto de \(f'(x)\) en \([-1, 1]\) ocurre en \(x = -1\): 
        \[ |f'(-1)| = |2(-1 - 0.5)| = 3 \quad \Rightarrow \quad l = 3 \]
    \end{itemize}

    \textbf{2. Fórmulas de intersección (Límite Inferior):}
    \begin{itemize}
        \item \textbf{Límites exteriores:}
        \[ y_{\text{izq}} = y_0 - l(x_0 - a) \quad \text{y} \quad y_{\text{der}} = y_{\text{fin}} - l(b - x_{\text{fin}}) \]
        \item \textbf{Puntos interiores} (entre nodos consecutivos \(A\) y \(B\)):
        \[ t = \frac{(A_y - B_y) - l(A_x - B_x)}{2l}, P_x = A_x + t, \quad P_y = A_y - l \cdot t \]
    \end{itemize}

    \textbf{3. Inicialización:}
    \(x_1 = \frac{-1 + 1}{2} = 0 \quad \Rightarrow \quad f(0) = 0.25\). Nodos iniciales: \(\{(0, 0.25)\}\).
\end{frame}

% 7. Paso a paso

\begin{frame}{Shubert-Piyavskii: Paso a paso}
\small
En cada iteración calculamos los posibles mínimos de la envolvente diente de sierra (\(y_{\text{izq}}\), \(y_{\text{der}}\) y los \(P_y\) interiores). Seleccionamos el menor (\(y_{\min}\)) en la posición \(x_{\min}\), evaluamos \(f(x_{\min})\) y actualizamos los puntos de soporte.

\vspace{0.25cm}
\centering
\renewcommand{\arraystretch}{1.25}
\resizebox{\textwidth}{!}{% 
\begin{tabular}{@{}c|l|cc|cc@{}}
    \toprule
    \textbf{Iter} & \textbf{Nodos Activos \((x, y)\)} & \textbf{Mín Estimado} & \(\bm{x_{\min}}\) & \(\bm{f(x_{\min})}\) & \(\bm{\Delta}\) \\
    \midrule
    \textbf{1} & \((0, 0.25)\) 
    & \(y_{\text{izq}}, y_{\text{der}} = -2.75\) & \(-1.00\) & \(2.25\) & \(5.00\) \\
    
    \textbf{2} & \((-1, 2.25), (0, 0.25)\) 
    & \(y_{\text{der}} = -2.75\) & \(1.00\) & \(0.25\) & \(3.00\) \\
    
    \textbf{3} & \((-1, 2.25), (0, 0.25), (1, 0.25)\) 
    & \(P_{y2} = -1.25\) & \(0.50\) & \(0.00\) & \(1.25\) \\
    
    \textbf{4} & \scriptsize \((-1, 2.25), (0, 0.25), (0.5, 0), (1, 0.25)\) 
    & \(P_{y2} = -0.625\) & \(0.2917\) & \(0.0434\) & \(0.6684\) \\
    \bottomrule
\end{tabular}%  
}

\vspace{0.25cm}
\raggedright
\footnotesize
\textbf{Detalle de la Iteración 4:} \\
El mínimo estimado proviene de la intersección entre \((0, 0.25)\) y \((0.5, 0)\): \\
\(t = \frac{(0.25 - 0.0) - 3(0 - 0.5)}{6} = \frac{1.75}{6} \approx 0.2917\), \quad 
\(P_x = 0 + 0.2917 = \mathbf{0.2917}\), \\
\(P_y = 0.25 - 3(0.2917) = \mathbf{-0.6250}\), \quad 
\(\Delta = f(0.2917) - (-0.6250) = 0.0434 + 0.6250 = \mathbf{0.6684}\).
\end{frame}


% ---------------------------------------------------------------
% Shubert-Piyavskii: Iteración 1 (Corregido visualmente)
% ---------------------------------------------------------------
\begin{frame}{Shubert-Piyavskii: Iteración 1}
    \vspace{-0.2cm}
    \begin{exampleblock}{Estado Inicial e Iteración 1}
        \footnotesize
        Con el nodo central \((0, 0.25)\), los conos de pendiente \(L=3\) generan en las fronteras \(y_{\text{izq}} = y_{\text{der}} = -2.75\). Por orden de exploración, se evalúa la frontera izquierda \(x^{(1)} = -1.00\) con \(f(-1.0) = 2.25\), trazando el nuevo rayo descendente hacia el interior.
    \end{exampleblock}

    \vspace{-0.1cm}
    \begin{center}
    \begin{tikzpicture}[
        x=3.8cm, y=0.7cm, font=\scriptsize,
        declare function={
            f(\x) = (\x - 0.5)^2;
        }
    ]
        % Marco exterior ajustado a las dimensiones del slide
        \draw[thick, GrisTexto] (-1.15, -3.1) rectangle (1.15, 2.7);

        % Eje horizontal y = 0
        \draw[GrisTexto, dashed, thin] (-1.15, 0) -- (1.15, 0);

        % Marcas en el eje X (Dibujadas hacia AFUERA del marco)
        \draw[GrisTexto] (-1.0, -3.1) -- (-1.0, -3.25) node[below, text=GrisTexto] {\(-1\)};
        \draw[GrisTexto] (0.0, -3.1) -- (0.0, -3.25) node[below, text=GrisTexto] {\(0\)};
        \draw[GrisTexto] (1.0, -3.1) -- (1.0, -3.25) node[below, text=GrisTexto] {\(1\)};

        % Marcas en el eje Y (Dibujadas hacia AFUERA del marco)
        \draw[GrisTexto] (-1.15, -2.75) -- (-1.25, -2.75) node[left, text=GrisTexto] {\(-2.75\)};
        \draw[GrisTexto] (-1.15, 0.0) -- (-1.25, 0.0) node[left, text=GrisTexto] {\(0\)};
        \draw[GrisTexto] (-1.15, 2.25) -- (-1.25, 2.25) node[left, text=GrisTexto] {\(2.25\)};

        % Leyenda distribuida horizontalmente para evitar superposición
        \draw[white, thick] (-1.0, 2.45) -- (-0.8, 2.45) node[right, text=white] {\(f(x) = (x-0.5)^2\)};
        \draw[NaranjaBase, thick] (-0.2, 2.45) -- (0.0, 2.45) node[right, text=NaranjaBase] {Envolvente inicial};
        \draw[CelesteBrillante, thick] (0.65, 2.45) -- (0.8, 2.45) node[right, text=CelesteBrillante] {Nuevo cono};

        % Función objetivo f(x)
        \draw[white, thick, domain=-1:1, samples=60, smooth] plot (\x, {f(\x)});

        % Envolvente inicial trazada desde (0, 0.25)
        \draw[NaranjaBase, thick] (-1.0, -2.75) -- (0.0, 0.25) -- (1.0, -2.75);

        % Punto inferior evaluado en x = -1
        \filldraw[NaranjaBase] (-1.0, -2.75) circle (1.5pt) node[below right, text=NaranjaBase, yshift=-1pt] {\((-1, -2.75)\)};

        % Flecha de evaluación vertical
        \draw[-stealth, dashed, white, thick] (-1.0, -2.6) -- (-1.0, 2.15);

        % Nuevo rayo exacto: desde (-1, 2.25) hasta la intersección (-0.1667, -0.25)
        \draw[CelesteBrillante, very thick] (-1.0, 2.25) -- (-0.1667, -0.25);
        \filldraw[CelesteBrillante] (-0.1667, -0.25) circle (1.3pt) node[below, text=CelesteBrillante, yshift=-3pt] {\scriptsize\((-0.17, -0.25)\)};

        % Nodos evaluados sobre f(x)
        \filldraw[CelesteBrillante] (0.0, 0.25) circle (1.6pt) node[above, text=CelesteBrillante, yshift=2pt] {\((0, 0.25)\)};
        % Reubicado para que no choque con el eje Y
        \filldraw[CelesteBrillante] (-1.0, 2.25) circle (1.6pt) node[above right, text=CelesteBrillante] {\(x^{(1)} = -1\)};
    \end{tikzpicture}
    \end{center}
\end{frame}


% ---------------------------------------------------------------
% Shubert-Piyavskii: Iteración 2
% ---------------------------------------------------------------
\begin{frame}{Shubert-Piyavskii: Iteración 2}
    \vspace{-0.2cm}
    \begin{exampleblock}{Iteración 2: Exploración de la frontera derecha}
        \footnotesize
        Tras actualizar la envolvente, el mínimo estimado (el vértice más profundo) se ubica en la frontera derecha en \(x^{(2)} = 1.00\), donde \(y_{\text{der}} = -2.75\). Al evaluar, obtenemos \(f(1.0) = 0.25\). Se traza un nuevo rayo hacia el interior que levanta la cota inferior por la derecha.
    \end{exampleblock}

    \vspace{-0.1cm}
    \begin{center}
    \begin{tikzpicture}[
        x=3.8cm, y=0.7cm, font=\scriptsize,
        declare function={
            f(\x) = (\x - 0.5)^2;
        }
    ]
        % Marco exterior ajustado a las dimensiones del slide
        \draw[thick, GrisTexto] (-1.15, -3.1) rectangle (1.15, 2.7);

        % Eje horizontal y = 0
        \draw[GrisTexto, dashed, thin] (-1.15, 0) -- (1.15, 0);

        % Marcas en el eje X (Dibujadas hacia AFUERA del marco)
        \draw[GrisTexto] (-1.0, -3.1) -- (-1.0, -3.25) node[below, text=GrisTexto] {\(-1\)};
        \draw[GrisTexto] (0.0, -3.1) -- (0.0, -3.25) node[below, text=GrisTexto] {\(0\)};
        \draw[GrisTexto] (1.0, -3.1) -- (1.0, -3.25) node[below, text=GrisTexto] {\(1\)};

        % Marcas en el eje Y (Dibujadas hacia AFUERA del marco)
        \draw[GrisTexto] (-1.15, -2.75) -- (-1.25, -2.75) node[left, text=GrisTexto] {\(-2.75\)};
        \draw[GrisTexto] (-1.15, 0.0) -- (-1.25, 0.0) node[left, text=GrisTexto] {\(0\)};
        \draw[GrisTexto] (-1.15, 2.25) -- (-1.25, 2.25) node[left, text=GrisTexto] {\(2.25\)};

        % Leyenda distribuida horizontalmente para evitar superposición
        \draw[white, thick] (-1.0, 2.45) -- (-0.8, 2.45) node[right, text=white] {\(f(x) = (x-0.5)^2\)};
        \draw[NaranjaBase, thick] (-0.2, 2.45) -- (0.0, 2.45) node[right, text=NaranjaBase] {Envolvente actual};
        \draw[CelesteBrillante, thick] (0.65, 2.45) -- (0.8, 2.45) node[right, text=CelesteBrillante] {Nuevo cono};

        % Función objetivo f(x)
        \draw[white, thick, domain=-1:1, samples=60, smooth] plot (\x, {f(\x)});

        % Envolvente actual consolidada (heredada de la iteración 1)
        % Recorrido: (-1, 2.25) -> intersección izq -> (0, 0.25) -> fondo derecho
        \draw[NaranjaBase, thick] (-1.0, 2.25) -- (-0.1667, -0.25) -- (0.0, 0.25) -- (1.0, -2.75);
        
        % Vértice consolidado de la iteración anterior (pasa a ser naranja)
        \filldraw[NaranjaBase] (-0.1667, -0.25) circle (1.3pt) node[below, text=NaranjaBase, yshift=-3pt] {\scriptsize\((-0.17, -0.25)\)};

        % Punto inferior evaluado en el paso actual: x = 1
        \filldraw[NaranjaBase] (1.0, -2.75) circle (1.5pt) node[below left, text=NaranjaBase, yshift=-1pt] {\((1, -2.75)\)};

        % Flecha de evaluación vertical (desde el mínimo inferior hasta f(1))
        \draw[-stealth, dashed, white, thick] (1.0, -2.6) -- (1.0, 0.15);

        % Nuevo rayo exacto: desde (1, 0.25) hacia la izquierda hasta la intersección (0.5, -1.25)
        \draw[CelesteBrillante, very thick] (1.0, 0.25) -- (0.5, -1.25);
        \filldraw[CelesteBrillante] (0.5, -1.25) circle (1.3pt) node[below, text=CelesteBrillante, yshift=-3pt] {\scriptsize\((0.5, -1.25)\)};

        % Nodos evaluados sobre f(x)
        \filldraw[CelesteBrillante] (-1.0, 2.25) circle (1.6pt); % Nodo antiguo
        \filldraw[CelesteBrillante] (0.0, 0.25) circle (1.6pt) node[above, text=CelesteBrillante, yshift=2pt] {\((0, 0.25)\)}; % Nodo central
        
        % Nodo recién evaluado, etiqueta desplazada a la izquierda para evitar el marco
        \filldraw[CelesteBrillante] (1.0, 0.25) circle (1.6pt) node[above left, text=CelesteBrillante] {\(x^{(2)} = 1\)};
    \end{tikzpicture}
    \end{center}
\end{frame}


% ---------------------------------------------------------------
% Shubert-Piyavskii: Iteración 3
% ---------------------------------------------------------------
\begin{frame}{Shubert-Piyavskii: Iteración 3}
    \vspace{-0.2cm}
    \begin{exampleblock}{Iteración 3: Exploración del interior}
        \footnotesize
        El mínimo estimado de la envolvente se encuentra ahora en el interior, en \(x^{(3)} = 0.50\), donde \(y_{\min} = -1.25\). Al evaluar, obtenemos \(f(0.5) = 0.00\) (el mínimo exacto). Se traza un nuevo cono desde este punto, levantando la envolvente y generando dos intersecciones simétricas.
    \end{exampleblock}

    \vspace{-0.1cm}
    \begin{center}
    \begin{tikzpicture}[
        x=3.8cm, y=0.7cm, font=\scriptsize,
        declare function={
            f(\x) = (\x - 0.5)^2;
        }
    ]
        % Marco exterior ajustado a las dimensiones del slide
        \draw[thick, GrisTexto] (-1.15, -3.1) rectangle (1.15, 2.7);

        % Eje horizontal y = 0
        \draw[GrisTexto, dashed, thin] (-1.15, 0) -- (1.15, 0);

        % Marcas en el eje X (Dibujadas hacia AFUERA del marco)
        \draw[GrisTexto] (-1.0, -3.1) -- (-1.0, -3.25) node[below, text=GrisTexto] {\(-1\)};
        \draw[GrisTexto] (0.0, -3.1) -- (0.0, -3.25) node[below, text=GrisTexto] {\(0\)};
        \draw[GrisTexto] (1.0, -3.1) -- (1.0, -3.25) node[below, text=GrisTexto] {\(1\)};

        % Marcas en el eje Y (Dibujadas hacia AFUERA del marco)
        \draw[GrisTexto] (-1.15, -2.75) -- (-1.25, -2.75) node[left, text=GrisTexto] {\(-2.75\)};
        \draw[GrisTexto] (-1.15, -1.25) -- (-1.25, -1.25) node[left, text=GrisTexto] {\(-1.25\)};
        \draw[GrisTexto] (-1.15, 0.0) -- (-1.25, 0.0) node[left, text=GrisTexto] {\(0\)};
        \draw[GrisTexto] (-1.15, 2.25) -- (-1.25, 2.25) node[left, text=GrisTexto] {\(2.25\)};

        % Leyenda distribuida horizontalmente
        \draw[white, thick] (-1.0, 2.45) -- (-0.8, 2.45) node[right, text=white] {\(f(x) = (x-0.5)^2\)};
        \draw[NaranjaBase, thick] (-0.2, 2.45) -- (0.0, 2.45) node[right, text=NaranjaBase] {Envolvente actual};
        \draw[CelesteBrillante, thick] (0.65, 2.45) -- (0.8, 2.45) node[right, text=CelesteBrillante] {Nuevo cono};

        % Función objetivo f(x)
        \draw[white, thick, domain=-1:1, samples=60, smooth] plot (\x, {f(\x)});

        % Envolvente actual consolidada (heredada de la iteración 2)
        \draw[NaranjaBase, thick] (-1.0, 2.25) -- (-0.1667, -0.25) -- (0.0, 0.25) -- (0.5, -1.25) -- (1.0, 0.25);
        
        % Vértice izquierdo consolidado
        \filldraw[NaranjaBase] (-0.1667, -0.25) circle (1.3pt);

        % Punto inferior evaluado en el paso actual: x = 0.5
        \filldraw[NaranjaBase] (0.5, -1.25) circle (1.5pt) node[below, text=NaranjaBase, yshift=-1pt] {\((0.5, -1.25)\)};

        % Flecha de evaluación vertical
        \draw[-stealth, dashed, white, thick] (0.5, -1.1) -- (0.5, -0.15);

        % Nuevo cono exacto: desde (0.5, 0.0) hacia izquierda y derecha
        % Intersección izquierda: (0.2917, -0.625)
        % Intersección derecha: (0.7083, -0.625)
        \draw[CelesteBrillante, very thick] (0.5, 0.0) -- (0.2917, -0.625);
        \draw[CelesteBrillante, very thick] (0.5, 0.0) -- (0.7083, -0.625);
        
        % Etiquetas de las nuevas intersecciones
        \filldraw[CelesteBrillante] (0.2917, -0.625) circle (1.3pt) node[below left, text=CelesteBrillante, xshift=6pt, yshift=-2pt] {\scriptsize\((0.29, -0.63)\)};
        \filldraw[CelesteBrillante] (0.7083, -0.625) circle (1.3pt) node[below right, text=CelesteBrillante, xshift=-6pt, yshift=-2pt] {\scriptsize\((0.71, -0.63)\)};

        % Nodos evaluados sobre f(x)
        \filldraw[CelesteBrillante] (-1.0, 2.25) circle (1.6pt);
        \filldraw[CelesteBrillante] (0.0, 0.25) circle (1.6pt) node[above, text=CelesteBrillante, yshift=2pt] {\((0, 0.25)\)};
        \filldraw[CelesteBrillante] (1.0, 0.25) circle (1.6pt);
        
        % Nodo recién evaluado
        \filldraw[CelesteBrillante] (0.5, 0.0) circle (1.6pt) node[above, text=CelesteBrillante, yshift=2pt] {\(x^{(3)} = 0.5\)};
    \end{tikzpicture}
    \end{center}
\end{frame}


% Implementación
\begin{frame}[fragile]{Shubert-Piyavskii: Implementación}
    \begin{lstlisting}[basicstyle=\ttfamily\tiny]
class Punto:
    def __init__(self, x, y):
        self.x = x
        self.y = y
def shubert_piyavskii(f, a, b, l, epsilon):
    pts = [Punto((a + b) / 2.0, f((a + b) / 2.0))]
    delta = float('inf')
    while delta > epsilon:
        best_i, best_x, best_y = 0, 0.0, float('inf')
        y_left = pts[0].y - l * (pts[0].x - a)
        if y_left < best_y: best_i, best_x, best_y = 0, a, y_left
        y_right = pts[-1].y - l * (b - pts[-1].x)
        if y_right < best_y: best_i, best_x, best_y = len(pts), b, y_right
        for i in range(1, len(pts)):
            A, B = pts[i-1], pts[i]
            t = ((A.y - B.y) - l * (A.x - B.x)) / (2 * l)
            P_x, P_y = A.x + t, A.y - l * t
            if P_y < best_y: best_i, best_x, best_y = i, P_x, P_y
        f_best_x = f(best_x)
        pts.insert(best_i, Punto(best_x, f_best_x))
        delta = f_best_x - best_y
    mejor_punto = min(pts, key=lambda p: p.y)
    return mejor_punto.x, mejor_punto.y
    \end{lstlisting}
\end{frame}


% ---------------------------------------------------------------
% 8. Análisis de Convergencia
% ---------------------------------------------------------------
\begin{frame}{Shubert-Piyavskii: Análisis de Convergencia}
    \begin{block}{Garantía de Óptimo Global}
        \footnotesize
        Shubert-Piyavskii garantiza la convergencia al \textbf{mínimo global}.
    \end{block}

    \vspace{0.2cm}
    \footnotesize
    \textbf{¿Por qué es capaz de encontrar el mínimo global?}
    \scriptsize
    \begin{itemize}
        \item \textbf{Límite Inferior:} La envolvente en forma de dientes de sierra es un límite inferior absoluto. Gracias a la constante de Lipschitz \(\ell\), sabemos que la función real \textit{nunca} descenderá por debajo.
        \item \textbf{Eliminación de Regiones:} Al levantar el límite inferior iterativamente, el algoritmo descarta áreas del dominio donde es matemáticamente imposible que exista un valor menor al mínimo actual encontrado.
        \item \textbf{Exploración del Peor:} Al seleccionar siempre el punto con el menor valor del límite inferior (\(y^*\)), el método ataca directamente la región de mayor ``incertidumbre'' o potencial, forzando a que la brecha \(\Delta = f(x^*) - y^*\) se reduzca.
    \end{itemize}

    \begin{alertblock}{Conclusión de Parada}
        \scriptsize
        Cuando la diferencia máxima entre el límite inferior y la función cae por debajo de \(\epsilon\) (\(f(x^*) - y^* < \epsilon\)), el método se detiene. En este punto \textbf{ningún otro punto en todo el dominio} puede tener un valor menor a \(f(x^*) - \epsilon\).
    \end{alertblock}
\end{frame}